\section{帝国能力：不是强弱刻度，而是一组连接}
\label{sec:synthesis-capacity}

汉帝国最有力的时刻，未必是所有命令都从长安或雒阳笔直抵达的时刻，而是中枢能让不同地区的命令、粮食、人员和信息暂时接续的时刻。早汉以关中仓储、宗室封国和功臣网络开始；削藩以后，选官、钱币、转运和郡县文书扩大了可被中枢看见的范围；东汉复兴则需要重新取得道路、城邑与簿籍。它们都不是一条“集权上升”的直线。每多一个接口，就多一处必须依赖地方中介、家族、将领或财政资源的地方。

这个视角使辅政、外戚、宦官和州牧不再只是道德化的反面人物。他们占据的正是皇帝难以亲自覆盖的节点：禁军、尚书、任命、边军与地方信息。节点在危机中能让制度继续运转，也会让占据它的人获得超出名义职位的力量。汉末并非国家能力突然归零，而是这些节点被不同的区域集团分别接管，不能再回到同一套优先次序中。

\subsection{制度得失}

以国家能力衡量，郡县文书、选官与财政调度提高了跨区域协作的可能；以社会成本衡量，它们也使征发、迁徙和竞争更深入地进入家庭与地方。王国、边部与地方网络并非只有阻碍作用，它们提供过中枢无法独自供给的兵员、信息与协商渠道。问题不在于有没有地方中介，而在于帝国是否仍能把中介的资源和责任接到一个共同的、可被检验的秩序中。
